\section{Discussion and Conclusion}
\label{sec:discussion}

The results support a specific claim: piecewise language is useful for both conventional and structured action heads, and an event-aligned spline head benefits more. We attribute this interaction to a better-matched interface. A clause specifies one manipulation phase; eight control points express its geometry; a gripper curve exposes its terminal interaction; and duration supplies its physical clock. The representation does not solve high-level planning by itself, but it reduces the burden placed on a single fixed-rate action array.

The language and timing results also explain one another. Words affect behavior only when demonstrations make them informative. Executable clauses vary within a long task and produce large gains; speed adverbs label physically different continuations and yield bidirectional pace control. By contrast, changing an object word without counterfactually varying the demonstrated target does not overcome a positional shortcut. Similarly, decode-time acceleration helps when demonstrations contain removable dead time (CALVIN), is neutral near actuator limits (RoboCasa), and trades accuracy for speed when scripted trajectories contain little slack (LIBERO). These conditions make the control surface predictable rather than a universal speedup claim.

Our study is presently limited to a 450M-parameter VLA, two claim-bearing LIBERO-Long tasks, and two RoboCasa programs. The order probe establishes causal first-subgoal steering but not robust reversed-order completion. These limitations identify the most direct next test: train on several subgoal combinations and hold out one composition, then evaluate a matched later-clause intervention. The real-robot protocol in Sec.~\ref{sec:hardware} instantiates the same test with drawer and stacking programs. Scaling \method{} to pretrained $\pi_{0.5}$-class backbones is complementary rather than necessary to identify the head-level effect.

In summary, language-addressable spline action programs connect three structures that fixed waypoint chunks leave entangled: semantic subgoals, continuous trajectory shape, and physical time. They preserve broad VLA competence, amplify executable-clause gains, causally steer subgoal order, and expose practical replanning, pace, and control-rate mechanisms from one checkpoint. This combination turns a compact spline from an output compression device into an interface between language and closed-loop robot execution.
